\documentclass[a4paper]{article}

\usepackage[a4paper,margin=2.2cm]{geometry}
\usepackage[T1]{fontenc}
\usepackage[utf8]{inputenc}
\usepackage[ngerman,english]{babel}
\usepackage{lmodern}
\usepackage{graphicx}
\usepackage{amsmath,amssymb}
\usepackage{minted}
\usepackage{multirow}
\usepackage{float}
\usepackage{todonotes}
\usepackage{forest}
\usepackage{xstring}
\usepackage{xcolor}
\usepackage[most]{tcolorbox}
\usepackage{xparse}
\usepackage{natbib}
\usepackage{hyperref}
\usepackage{enumitem}
\usepackage{parskip}

\usetikzlibrary{graphs,graphdrawing}
\usegdlibrary{layered}

% for external graphviz
\input{/home/donkarlo/Dropbox/repo/nd_ai_project/src/nd_ai/knoweldge_representation/language/formal/computer/programming/tex/latex/code_clips/external_graphviz.tex}

% for tags
\input{/home/donkarlo/Dropbox/repo/nd_ai_project/src/nd_ai/knoweldge_representation/language/formal/computer/programming/tex/latex/parts/control_sequence/kind/semtag/semtag.tex}

% for links
\input{/home/donkarlo/Dropbox/repo/nd_ai_project/src/nd_ai/knoweldge_representation/language/formal/computer/programming/tex/latex/code_clips/seehere.tex}

\title{Grammatikthemen im Deutschen}
\date{}

\begin{document}
    \maketitle
    
    
    \section{Verben {\small (Verbs)}}
        
        \subsection{Verbformen und Konjugation {\small (Verb Forms and Conjugation)}}
            
            \paragraph{Konjugation im Präsens {\small (Conjugation in the Present Tense)}}
                Im Präsens wird das Verb an Person und Numerus angepasst. Diese Form wird für Gegenwart, allgemeine Aussagen und manchmal auch für zukünftige Handlungen verwendet.\\
                In the present tense, the verb changes according to person and number. This form is used for the present, general statements, and sometimes also for future actions.
                
                \textbf{Beispiele:}\\
                Ich lerne jeden Tag Deutsch. : I learn German every day.\\
                Du arbeitest heute sehr schnell. : You are working very quickly today.\\
                Er wohnt seit zwei Jahren in Wien. : He has been living in Vienna for two years.\\
                Wir fahren morgen nach Graz. : We are going to Graz tomorrow.
            
            \paragraph{Starke und schwache Verben {\small (Strong and Weak Verbs)}}
                Schwache Verben bilden ihre Formen regelmäßig, während starke Verben oft den Stammvokal ändern. Dieser Unterschied ist besonders in den Vergangenheitsformen wichtig.\\
                Weak verbs form their tenses regularly, while strong verbs often change their stem vowel. This distinction is especially important in past forms.
                
                \textbf{Beispiele:}\\
                Ich machte die Tür zu. : I closed the door.\\
                Er kaufte ein neues Buch. : He bought a new book.\\
                Sie ging nach Hause. : She went home.\\
                Wir fanden die Lösung schnell. : We found the solution quickly.
            
            \paragraph{Gemischte Verben {\small (Mixed Verbs)}}
                Gemischte Verben verbinden Merkmale starker und schwacher Verben. Sie ändern oft den Stammvokal, haben aber schwache Endungen.\\
                Mixed verbs combine features of strong and weak verbs. They often change the stem vowel but still take weak endings.
                
                \textbf{Beispiele:}\\
                Ich dachte lange darüber nach. : I thought about it for a long time.\\
                Er brachte mir die Unterlagen. : He brought me the documents.\\
                Sie kannte die Antwort schon. : She already knew the answer.\\
                Wir nannten ihn beim Vornamen. : We called him by his first name.
            
            \paragraph{Trennbare und nicht trennbare Verben {\small (Separable and Inseparable Verbs)}}
                Trennbare Verben teilen sich im Hauptsatz, während nicht trennbare Verben zusammenbleiben. Die Bedeutung ändert sich oft mit dem Präfix.\\
                Separable verbs split in main clauses, whereas inseparable verbs remain together. The meaning often changes with the prefix.
                
                \textbf{Beispiele:}\\
                Ich stehe jeden Morgen um sechs Uhr auf. : I get up at six every morning.\\
                Er ruft seine Mutter später an. : He calls his mother later.\\
                Sie versteht die Frage nicht. : She does not understand the question.\\
                Wir besuchen unsere Freunde am Wochenende. : We visit our friends on the weekend.
            
            \paragraph{Reflexive Verben {\small (Reflexive Verbs)}}
                Reflexive Verben stehen mit einem Reflexivpronomen. Manche Verben sind immer reflexiv, andere nur in bestimmten Bedeutungen.\\
                Reflexive verbs are used with a reflexive pronoun. Some verbs are always reflexive, while others are reflexive only in certain meanings.
                
                \textbf{Beispiele:}\\
                Ich interessiere mich für Geschichte. : I am interested in history.\\
                Du erinnerst dich an den Termin. : You remember the appointment.\\
                Er setzt sich auf den Stuhl. : He sits down on the chair.\\
                Wir freuen uns auf den Urlaub. : We are looking forward to the vacation.
            
            \paragraph{Modalverben {\small (Modal Verbs)}}
                Modalverben verändern die Bedeutung eines Vollverbs und drücken Möglichkeit, Notwendigkeit, Erlaubnis, Wunsch oder Pflicht aus.\\
                Modal verbs modify the meaning of a main verb and express possibility, necessity, permission, desire, or obligation.
                
                \textbf{Beispiele:}\\
                Ich muss heute länger arbeiten. : I have to work longer today.\\
                Du darfst hier nicht rauchen. : You may not smoke here.\\
                Er kann sehr gut schwimmen. : He can swim very well.\\
                Wir wollen bald umziehen. : We want to move soon.
    
    \subsection{Zeiten {\small (Tenses)}}
        
        \paragraph{Präsens {\small (Present)}}
            Das Präsens beschreibt Gegenwart, Gewohnheiten, allgemeine Tatsachen und in vielen Fällen auch Zukunft.\\
            The present tense describes the present, habits, general facts, and in many cases also the future.
            
            \textbf{Beispiele:}\\
            Ich lese gerade einen Artikel. : I am reading an article right now.\\
            Sie arbeitet im Krankenhaus. : She works in a hospital.\\
            Wasser kocht bei hundert Grad. : Water boils at one hundred degrees.\\
            Morgen fahre ich nach Linz. : Tomorrow I am going to Linz.
        
        \paragraph{Perfekt {\small (Perfect)}}
            Das Perfekt wird in der gesprochenen Sprache sehr oft für Vergangenes verwendet. Es wird mit \textit{haben} oder \textit{sein} und dem Partizip II gebildet.\\
            The perfect tense is very often used in spoken language for past events. It is formed with \textit{haben} or \textit{sein} and the past participle.
            
            \textbf{Beispiele:}\\
            Ich habe den Brief geschrieben. : I have written the letter.\\
            Er hat den Schlüssel verloren. : He has lost the key.\\
            Sie ist sehr früh angekommen. : She arrived very early.\\
            Wir haben lange auf dich gewartet. : We waited for you for a long time.
        
        \paragraph{Präteritum {\small (Simple Past)}}
            Das Präteritum wird häufig in schriftlichen Texten und bei bestimmten Verben wie \textit{sein}, \textit{haben} und den Modalverben verwendet.\\
            The simple past is frequently used in written texts and with certain verbs such as \textit{sein}, \textit{haben}, and the modal verbs.
            
            \textbf{Beispiele:}\\
            Ich war gestern zu Hause. : I was at home yesterday.\\
            Er hatte keine Zeit. : He had no time.\\
            Sie konnte die Aufgabe nicht lösen. : She could not solve the task.\\
            Wir gingen am Abend spazieren. : We went for a walk in the evening.
        
        \paragraph{Plusquamperfekt {\small (Past Perfect)}}
            Das Plusquamperfekt beschreibt eine Vorvergangenheit, also etwas, das vor einem anderen vergangenen Ereignis passiert war.\\
            The past perfect describes an action that happened before another action in the past.
            
            \textbf{Beispiele:}\\
            Ich hatte schon gegessen, bevor er kam. : I had already eaten before he came.\\
            Sie war gegangen, als ich anrief. : She had left when I called.\\
            Wir hatten die Arbeit beendet, bevor der Chef kam. : We had finished the work before the boss arrived.\\
            Er hatte das Buch gelesen, bevor der Film erschien. : He had read the book before the film came out.
        
        \paragraph{Futur I {\small (Future I)}}
            Das Futur I wird für Zukunft, Vermutung oder Absicht verwendet. Oft wird aber auch Präsens mit Zeitangabe benutzt.\\
            Future I is used for future actions, assumptions, or intentions. However, the present tense with a time expression is also often used.
            
            \textbf{Beispiele:}\\
            Ich werde morgen anrufen. : I will call tomorrow.\\
            Er wird später kommen. : He will come later.\\
            Sie wird das schon wissen. : She probably knows that already.\\
            Wir werden bald fertig sein. : We will be finished soon.
        
        \paragraph{Futur II {\small (Future Perfect)}}
            Das Futur II drückt aus, dass etwas in der Zukunft abgeschlossen sein wird, oder es dient zur Vermutung über Vergangenes.\\
            Future II expresses that something will have been completed in the future, or it is used to make assumptions about the past.
            
            \textbf{Beispiele:}\\
            Ich werde die Arbeit bis morgen beendet haben. : I will have finished the work by tomorrow.\\
            Er wird schon angekommen sein. : He will probably have arrived already.\\
            Sie wird das Buch gelesen haben. : She will have read the book.\\
            Wir werden bis dahin umgezogen sein. : We will have moved by then.
    
    \subsection{Konjunktiv {\small (Subjunctive)}}
    
    \paragraph{Konjunktiv II {\small (Subjunctive II)}}
        Der Konjunktiv II wird für Irrealität, Wünsche, Höflichkeit und vorsichtige Aussagen verwendet.\\
        The subjunctive II is used for unreality, wishes, politeness, and cautious statements.
        
        \textbf{Beispiele:}\\
        Wenn ich mehr Zeit hätte, würde ich mehr lesen. : If I had more time, I would read more.\\
        Ich wäre jetzt gern am Meer. : I would like to be at the sea now.\\
        Könnten Sie mir bitte helfen? : Could you please help me?\\
        An deiner Stelle würde ich früher anfangen. : If I were you, I would start earlier.
    
    \paragraph{Konjunktiv II der Vergangenheit {\small (Past Subjunctive II)}}
        Diese Form drückt aus, dass etwas in der Vergangenheit anders hätte sein können oder sollen.\\
        This form expresses that something in the past could or should have been different.
        
        \textbf{Beispiele:}\\
        Ich hätte früher kommen sollen. : I should have come earlier.\\
        Er wäre fast gestürzt. : He almost would have fallen.\\
        Wir hätten das wissen müssen. : We should have known that.\\
        Sie hätte die Prüfung bestehen können. : She could have passed the exam.
    
    \paragraph{Konjunktiv I {\small (Subjunctive I)}}
        Der Konjunktiv I wird besonders in der indirekten Rede verwendet. Er schafft Distanz zum Gesagten und ist typisch für Berichte und Nachrichten.\\
        Subjunctive I is used especially in indirect speech. It creates distance from what is being said and is typical in reports and news.
        
        \textbf{Beispiele:}\\
        Er sagte, er sei krank. : He said that he was ill.\\
        Sie meinte, sie habe keine Zeit. : She said that she had no time.\\
        Der Zeuge erklärte, er habe nichts gesehen. : The witness stated that he had seen nothing.\\
        Die Zeitung berichtet, der Minister trete zurück. : The newspaper reports that the minister is resigning.
    
    \paragraph{Irreale Konditionalsätze {\small (Unreal Conditional Clauses)}}
        Irreale Konditionalsätze zeigen eine Bedingung, die nicht erfüllt ist oder nur vorgestellt wird.\\
        Unreal conditional clauses show a condition that is not fulfilled or only imagined.
        
        \textbf{Beispiele:}\\
        Wenn ich reich wäre, würde ich mehr reisen. : If I were rich, I would travel more.\\
        Wenn wir früher losgefahren wären, hätten wir den Zug erreicht. : If we had left earlier, we would have caught the train.\\
        Wenn sie besser Deutsch könnte, fände sie leichter Arbeit. : If she spoke better German, she would find work more easily.\\
        Wenn er mich gefragt hätte, hätte ich geholfen. : If he had asked me, I would have helped.
    
    \paragraph{Wunschsätze {\small (Wish Sentences)}}
        Wunschsätze drücken einen oft unerfüllbaren oder unwahrscheinlichen Wunsch aus.\\
        Wish sentences express a wish that is often unrealizable or unlikely.
        
        \textbf{Beispiele:}\\
        Wenn ich doch mehr Zeit hätte. : If only I had more time.\\
        Wäre ich jetzt nur zu Hause. : If only I were at home now.\\
        Hätte er mich doch angerufen. : If only he had called me.\\
        Wenn es morgen bloß nicht regnen würde. : If only it would not rain tomorrow.
    
    \subsection{Verbmuster und Ergänzungen {\small (Verb Patterns and Complements)}}
    
    \paragraph{Verben mit Akkusativobjekt {\small (Verbs with an Accusative Object)}}
        Viele Verben verlangen ein direktes Objekt im Akkusativ.\\
        Many verbs require a direct object in the accusative case.
        
        \textbf{Beispiele:}\\
        Ich lese das Buch. : I am reading the book.\\
        Sie kauft einen Mantel. : She is buying a coat.\\
        Er sucht seinen Schlüssel. : He is looking for his key.\\
        Wir brauchen eine Lösung. : We need a solution.
    
    \paragraph{Verben mit Dativobjekt {\small (Verbs with a Dative Object)}}
        Einige Verben stehen mit einem Objekt im Dativ.\\
        Some verbs take an object in the dative case.
        
        \textbf{Beispiele:}\\
        Ich helfe meinem Bruder. : I help my brother.\\
        Sie dankt ihrer Lehrerin. : She thanks her teacher.\\
        Er folgt dem Arzt. : He follows the doctor.\\
        Wir vertrauen dem System nicht. : We do not trust the system.
    
    \paragraph{Verben mit Dativ- und Akkusativobjekt {\small (Verbs with Dative and Accusative Objects)}}
        Bestimmte Verben verbinden zwei Objekte miteinander. Häufig steht eine Person im Dativ und eine Sache im Akkusativ.\\
        Certain verbs combine two objects. Often a person is in the dative and a thing is in the accusative.
        
        \textbf{Beispiele:}\\
        Ich gebe dem Kind ein Buch. : I give the child a book.\\
        Sie erklärt mir die Aufgabe. : She explains the task to me.\\
        Er schickt seiner Mutter eine Nachricht. : He sends his mother a message.\\
        Wir zeigen den Gästen die Stadt. : We show the guests the city.
    
    \paragraph{Verben mit Präpositionalobjekt {\small (Verbs with Prepositional Objects)}}
        Viele Verben stehen fest mit einer bestimmten Präposition.\\
        Many verbs are fixed with a specific preposition.
        
        \textbf{Beispiele:}\\
        Ich warte auf den Bus. : I am waiting for the bus.\\
        Sie interessiert sich für Kunst. : She is interested in art.\\
        Er denkt an seine Kindheit. : He thinks about his childhood.\\
        Wir sprechen über das Problem. : We are talking about the problem.
    
    \paragraph{Verben mit Pronominaladverbien {\small (Verbs with Pronominal Adverbs)}}
        Pronominaladverbien wie \textit{darauf}, \textit{damit}, \textit{woran} ersetzen Präpositionalobjekte mit Sachen.\\
        Pronominal adverbs such as \textit{darauf}, \textit{damit}, \textit{woran} replace prepositional objects referring to things.
        
        \textbf{Beispiele:}\\
        Ich warte darauf. : I am waiting for that.\\
        Sie spricht oft darüber. : She often talks about it.\\
        Woran denkst du? : What are you thinking about?\\
        Damit bin ich nicht einverstanden. : I do not agree with that.
    
    \paragraph{Verben mit festen Ergänzungsmustern {\small (Verbs with Fixed Complement Patterns)}}
        Manche Verben haben feste Strukturen, zum Beispiel mit Infinitiv, Nebensatz oder bestimmten Präpositionen.\\
        Some verbs have fixed structures, for example with an infinitive, a subordinate clause, or specific prepositions.
        
        \textbf{Beispiele:}\\
        Ich bitte dich, mir zu helfen. : I ask you to help me.\\
        Sie hofft, dass alles gut geht. : She hopes that everything goes well.\\
        Er beginnt, Deutsch zu lernen. : He begins to learn German.\\
        Wir hindern ihn daran, sofort zu gehen. : We prevent him from leaving immediately.
    
    \subsection{Verbale Spezialthemen {\small (Special Verb Topics)}}
    
    \paragraph{Bedeutung und Funktion von \textit{werden} {\small (Meaning and Function of \textit{werden})}}
        \textit{Werden} kann Zukunft, Veränderung oder Passiv ausdrücken. Es ist eines der zentralen Verben der deutschen Grammatik.\\
        \textit{Werden} can express future, change, or passive voice. It is one of the central verbs in German grammar.
        
        \textbf{Beispiele:}\\
        Ich werde morgen anrufen. : I will call tomorrow.\\
        Es wird kalt. : It is getting cold.\\
        Das Haus wird renoviert. : The house is being renovated.\\
        Er wird Arzt. : He becomes a doctor.
    
    \paragraph{\textit{brauchen} und \textit{sich lassen} {\small (\textit{brauchen} and \textit{sich lassen})}}
        \textit{Brauchen} mit \textit{zu} erscheint oft in Verneinung. \textit{Sich lassen} drückt Möglichkeit oder Machbarkeit aus.\\
        \textit{Brauchen} with \textit{zu} often appears in negation. \textit{Sich lassen} expresses possibility or feasibility.
        
        \textbf{Beispiele:}\\
        Du brauchst heute nicht zu kommen. : You do not need to come today.\\
        Man braucht das nicht extra zu erklären. : There is no need to explain that separately.\\
        Das Problem lässt sich lösen. : The problem can be solved.\\
        Die Tür lässt sich nicht öffnen. : The door cannot be opened.
    
    \paragraph{Nomen-Verb-Verbindungen {\small (Noun-Verb Combinations)}}
        Nomen-Verb-Verbindungen sind häufig und oft stilistisch wichtig. Manche sind frei, andere fest oder idiomatisch.\\
        Noun-verb combinations are common and often important stylistically. Some are free, others fixed or idiomatic.
        
        \textbf{Beispiele:}\\
        eine Entscheidung treffen : to make a decision\\
        in Frage kommen : to be considered\\
        Abschied nehmen : to say goodbye\\
        zur Verfügung stehen : to be available
    
    \paragraph{Funktionsverbgefüge {\small (Light Verb Constructions)}}
        Funktionsverbgefüge sind feste Verbindungen aus Verb und Nomen. Sie sind besonders häufig in formellen Texten.\\
        Light verb constructions are fixed combinations of a verb and a noun. They are especially common in formal texts.
        
        \textbf{Beispiele:}\\
        eine Rolle spielen : to play a role\\
        zur Anwendung kommen : to be applied\\
        eine Entscheidung treffen : to make a decision\\
        in Betracht ziehen : to take into consideration
    
    \paragraph{\textit{haben/sein} + \textit{zu} + Infinitiv {\small (\textit{have/to be} + infinitive)}}
        Diese Konstruktion drückt Notwendigkeit, Verpflichtung oder Möglichkeit aus.\\
        This construction expresses necessity, obligation, or possibility.
        
        \textbf{Beispiele:}\\
        Ich habe noch viel zu erledigen. : I still have a lot to do.\\
        Der Bericht ist bis Freitag zu schreiben. : The report is to be written by Friday.\\
        Das Problem ist nicht leicht zu lösen. : The problem is not easy to solve.\\
        Sie hat morgen einen Vortrag zu halten. : She has to give a presentation tomorrow.
    
    \paragraph{\textit{scheinen} + \textit{zu} + Infinitiv {\small (\textit{seem} + infinitive)}}
        Mit \textit{scheinen} formuliert man vorsichtige Einschätzungen oder Beobachtungen.\\
        With \textit{scheinen}, one formulates cautious assessments or observations.
        
        \textbf{Beispiele:}\\
        Er scheint müde zu sein. : He seems to be tired.\\
        Sie scheint die Antwort zu kennen. : She seems to know the answer.\\
        Das Problem scheint größer zu werden. : The problem seems to be getting bigger.\\
        Die Lage scheint sich zu beruhigen. : The situation seems to be calming down.
    
    
    \section{Nomen und Artikel {\small (Nouns and Articles)}}
    
    \subsection{Nomen {\small (Nouns)}}
        
        \paragraph{Bedeutung und Form der Nomen {\small (Meaning and Form of Nouns)}}
            Nomen bezeichnen Personen, Dinge, Begriffe oder Sachverhalte. Im Deutschen werden sie großgeschrieben.\\
            Nouns refer to people, things, concepts, or states of affairs. In German, they are capitalized.
            
            \textbf{Beispiele:}\\
            der Tisch : the table\\
            die Hoffnung : hope\\
            das Studium : studies\\
            die Entscheidung : decision
        
        \paragraph{Genus {\small (Gender)}}
            Jedes Nomen hat ein grammatisches Geschlecht: maskulin, feminin oder neutral. Das Genus muss oft mit dem Artikel gelernt werden.\\
            Every noun has a grammatical gender: masculine, feminine, or neuter. The gender often has to be learned together with the article.
            
            \textbf{Beispiele:}\\
            der Tag : the day\\
            die Sprache : the language\\
            das Problem : the problem\\
            der Gedanke : the thought
        
        \paragraph{Numerus {\small (Number)}}
            Nomen können im Singular oder im Plural stehen. Die Pluralbildung ist im Deutschen vielfältig und muss oft mitgelernt werden.\\
            Nouns can appear in the singular or plural. Plural formation in German is varied and often has to be learned individually.
            
            \textbf{Beispiele:}\\
            das Buch -- die Bücher : the book -- the books\\
            die Frau -- die Frauen : the woman -- the women\\
            der Student -- die Studenten : the student -- the students\\
            das Kind -- die Kinder : the child -- the children
        
        \paragraph{Kasus {\small (Case)}}
            Nomen verändern zusammen mit Artikeln und Adjektiven ihre Form je nach Funktion im Satz.\\
            Nouns change their form together with articles and adjectives according to their function in the sentence.
            
            \textbf{Beispiele:}\\
            Der Mann kommt. : The man is coming.\\
            Ich sehe den Mann. : I see the man.\\
            Ich helfe dem Mann. : I help the man.\\
            Das Auto des Mannes ist neu. : The man's car is new.
        
        \paragraph{n-Deklination {\small (n-Declension)}}
            Bestimmte maskuline Nomen bekommen in fast allen Fällen außer dem Nominativ Singular die Endung \textit{-n} oder \textit{-en}.\\
            Certain masculine nouns take the ending \textit{-n} or \textit{-en} in almost all cases except the nominative singular.
            
            \textbf{Beispiele:}\\
            der Student -- des Studenten : the student -- of the student\\
            der Kollege -- dem Kollegen : the colleague -- to the colleague\\
            Ich sehe den Jungen. : I see the boy.\\
            Er spricht mit dem Kunden. : He speaks with the customer.
        
        \paragraph{Komposita {\small (Compound Nouns)}}
            Im Deutschen lassen sich viele Nomen zu einem neuen Wort verbinden. Meistens bestimmt das letzte Element das Genus.\\
            In German, many nouns can be combined into a new word. Usually, the last element determines the gender.
            
            \textbf{Beispiele:}\\
            das Wörterbuch : dictionary\\
            die Fahrkarte : ticket\\
            der Sprachkurs : language course\\
            das Arbeitszimmer : study room
        
        \paragraph{Nominalisierung {\small (Nominalization)}}
            Verben und Adjektive können zu Nomen gemacht werden. Diese Form ist besonders häufig in formeller Sprache.\\
            Verbs and adjectives can be turned into nouns. This form is especially frequent in formal language.
            
            \textbf{Beispiele:}\\
            das Lesen : reading\\
            beim Schreiben : while writing\\
            das Gute : the good\\
            das Lernen fällt mir leicht. : Learning is easy for me.
    
    \subsection{Artikel {\small (Articles)}}
    
    \paragraph{Bestimmte Artikel {\small (Definite Articles)}}
        Der bestimmte Artikel verweist auf etwas Bekanntes oder konkret Gemeintes.\\
        The definite article refers to something known or specific.
        
        \textbf{Beispiele:}\\
        Der Lehrer ist krank. : The teacher is ill.\\
        Ich nehme den Bus. : I am taking the bus.\\
        Die Wohnung ist groß. : The apartment is big.\\
        Das Problem ist bekannt. : The problem is known.
    
    \paragraph{Unbestimmte Artikel {\small (Indefinite Articles)}}
        Der unbestimmte Artikel nennt etwas Neues oder Nichtgenaues.\\
        The indefinite article introduces something new or nonspecific.
        
        \textbf{Beispiele:}\\
        Ich suche eine Wohnung. : I am looking for an apartment.\\
        Er hat ein Auto gekauft. : He bought a car.\\
        Sie braucht einen Termin. : She needs an appointment.\\
        Wir haben einen Arzt gefunden. : We found a doctor.
    
    \paragraph{Nullartikel {\small (Zero Article)}}
        In bestimmten Fällen steht kein Artikel, besonders bei Stoffnamen, Berufen nach \textit{sein} oder bei abstrakten Begriffen.\\
        In certain cases, no article is used, especially with materials, professions after \textit{sein}, or abstract concepts.
        
        \textbf{Beispiele:}\\
        Er ist Arzt. : He is a doctor.\\
        Ich trinke Wasser. : I drink water.\\
        Liebe ist wichtig. : Love is important.\\
        Wir lernen Deutsch. : We are learning German.
    
    \paragraph{Negation mit \textit{kein} {\small (Negation with \textit{kein})}}
        \textit{Kein} verneint Nomen mit unbestimmtem Artikel oder ohne Artikel.\\
        \textit{Kein} negates nouns with an indefinite article or without an article.
        
        \textbf{Beispiele:}\\
        Ich habe kein Auto. : I do not have a car.\\
        Sie trinkt keinen Kaffee. : She does not drink coffee.\\
        Er hat keine Zeit. : He has no time.\\
        Wir machen keinen Fehler. : We are making no mistake.
    
    
    \section{Pronomen und Artikelwörter {\small (Pronouns and Determiners)}}
    
    \subsection{Personalpronomen {\small (Personal Pronouns)}}
        
        \paragraph{Personalpronomen im Nominativ, Akkusativ und Dativ {\small (Personal Pronouns in Nominative, Accusative, and Dative)}}
            Personalpronomen ersetzen Nomen und verändern ihre Form nach Kasus.\\
            Personal pronouns replace nouns and change their form according to case.
            
            \textbf{Beispiele:}\\
            Ich kenne ihn gut. : I know him well.\\
            Sie gibt mir das Buch. : She gives me the book.\\
            Wir sehen euch morgen. : We will see you tomorrow.\\
            Kannst du uns helfen? : Can you help us?
        
        \paragraph{Possessivartikel {\small (Possessive Determiners)}}
            Possessivartikel zeigen Besitz oder Zugehörigkeit an und passen sich an das Nomen an.\\
            Possessive determiners show possession or belonging and agree with the noun.
            
            \textbf{Beispiele:}\\
            Das ist mein Buch. : That is my book.\\
            Sie sucht ihre Tasche. : She is looking for her bag.\\
            Wir besuchen unsere Freunde. : We are visiting our friends.\\
            Er hat seinen Pass verloren. : He lost his passport.
        
        \paragraph{Demonstrativartikel und Demonstrativpronomen {\small (Demonstrative Determiners and Pronouns)}}
            Demonstrativformen heben etwas hervor oder verweisen auf etwas Bestimmtes.\\
            Demonstrative forms emphasize something or refer to something specific.
            
            \textbf{Beispiele:}\\
            Dieser Film ist interessant. : This film is interesting.\\
            Diese Frage ist wichtig. : This question is important.\\
            Das gefällt mir nicht. : I do not like that.\\
            Derjenige, der fragt, lernt. : The one who asks learns.
        
        \paragraph{Unbestimmte Pronomen und Artikel {\small (Indefinite Pronouns and Determiners)}}
            Diese Formen bezeichnen unbestimmte Personen, Dinge oder Mengen.\\
            These forms refer to indefinite persons, things, or quantities.
            
            \textbf{Beispiele:}\\
            Jemand wartet vor der Tür. : Someone is waiting at the door.\\
            Niemand wusste die Antwort. : Nobody knew the answer.\\
            Einige Studenten fehlen heute. : Some students are absent today.\\
            Alles war ruhig. : Everything was quiet.
        
        \paragraph{Pronominaladverbien {\small (Pronominal Adverbs)}}
            Pronominaladverbien ersetzen Präposition + Sache. Sie sind im Deutschen sehr häufig.\\
            Pronominal adverbs replace preposition + thing. They are very common in German.
            
            \textbf{Beispiele:}\\
            Daran denke ich oft. : I often think about that.\\
            Darüber sprechen wir morgen. : We will talk about that tomorrow.\\
            Womit schreibst du? : What are you writing with?\\
            Darauf freue ich mich. : I am looking forward to that.
        
        \paragraph{Relativpronomen {\small (Relative Pronouns)}}
            Relativpronomen leiten Relativsätze ein und beziehen sich auf ein Nomen im Hauptsatz.\\
            Relative pronouns introduce relative clauses and refer back to a noun in the main clause.
            
            \textbf{Beispiele:}\\
            Das ist die Frau, die hier arbeitet. : That is the woman who works here.\\
            Das ist der Mann, den ich gestern gesehen habe. : That is the man I saw yesterday.\\
            Das ist das Haus, in dem ich wohne. : That is the house in which I live.\\
            Das ist die Kollegin, deren Bruder Arzt ist. : That is the colleague whose brother is a doctor.
    
    
    \section{Adjektive und Steigerung {\small (Adjectives and Comparison)}}
    
    \subsection{Adjektive {\small (Adjectives)}}
        
        \paragraph{Adjektive und Adjektivdeklination {\small (Adjectives and Adjective Declension)}}
            Adjektive beschreiben Eigenschaften. Vor Nomen werden sie dekliniert, nach \textit{sein} oder \textit{werden} bleiben sie unverändert.\\
            Adjectives describe qualities. Before nouns they are declined; after \textit{sein} or \textit{werden} they remain unchanged.
            
            \textbf{Beispiele:}\\
            Das ist ein interessantes Buch. : That is an interesting book.\\
            Sie trägt einen langen Mantel. : She is wearing a long coat.\\
            Der Film war spannend. : The film was exciting.\\
            Wir wohnen in einer kleinen Wohnung. : We live in a small apartment.
        
        \paragraph{Komparativ {\small (Comparative)}}
            Der Komparativ vergleicht zwei Dinge oder Personen miteinander.\\
            The comparative compares two things or people.
            
            \textbf{Beispiele:}\\
            Dieses Buch ist interessanter als das andere. : This book is more interesting than the other one.\\
            Er arbeitet schneller als ich. : He works faster than I do.\\
            Die Aufgabe ist leichter als erwartet. : The task is easier than expected.\\
            Heute ist es wärmer als gestern. : Today it is warmer than yesterday.
        
        \paragraph{Superlativ {\small (Superlative)}}
            Der Superlativ bezeichnet den höchsten Grad einer Eigenschaft.\\
            The superlative expresses the highest degree of a quality.
            
            \textbf{Beispiele:}\\
            Das ist der beste Vorschlag. : That is the best suggestion.\\
            Sie ist am schnellsten gelaufen. : She ran the fastest.\\
            Dieser Weg ist der kürzeste. : This path is the shortest.\\
            Heute ist es am kältesten. : Today it is the coldest.
        
        \paragraph{Vergleichsformen mit \textit{so ... wie} und \textit{als} {\small (Comparison with \textit{so ... wie} and \textit{als})}}
            Mit \textit{so ... wie} vergleicht man Gleichheit, mit \textit{als} Ungleichheit.\\
            \textit{So ... wie} is used for equality, and \textit{als} for inequality.
            
            \textbf{Beispiele:}\\
            Er ist so groß wie sein Bruder. : He is as tall as his brother.\\
            Sie spricht besser Deutsch als ich. : She speaks German better than I do.\\
            Das Auto ist nicht so teuer wie gedacht. : The car is not as expensive as thought.\\
            Heute bin ich müder als gestern. : Today I am more tired than yesterday.
        
        \paragraph{Vergleich mit \textit{je ... desto/umso} {\small (Comparison with \textit{the ... the})}}
            Diese Struktur zeigt eine proportionale Entwicklung zwischen zwei Sachverhalten.\\
            This structure shows a proportional development between two situations.
            
            \textbf{Beispiele:}\\
            Je mehr du lernst, desto sicherer wirst du. : The more you study, the more confident you become.\\
            Je länger ich warte, umso ungeduldiger werde ich. : The longer I wait, the more impatient I become.\\
            Je früher wir losfahren, desto besser ist es. : The earlier we leave, the better it is.\\
            Je älter er wird, desto ruhiger wird er. : The older he gets, the calmer he becomes.
        
        \paragraph{Partizip I und Partizip II als Adjektive {\small (Present and Past Participles as Adjectives)}}
            Partizipien können wie Adjektive verwendet werden und sind besonders in schriftlicher Sprache häufig.\\
            Participles can be used like adjectives and are especially common in written language.
            
            \textbf{Beispiele:}\\
            das schlafende Kind : the sleeping child\\
            die steigenden Preise : the rising prices\\
            der geschriebene Text : the written text\\
            die geschlossene Tür : the closed door
        
        \paragraph{Erweiterte Partizipialattribute {\small (Extended Participial Attributes)}}
            Erweiterte Partizipialattribute verdichten Informationen und sind typisch für formellere Texte.\\
            Extended participial attributes condense information and are typical of more formal texts.
            
            \textbf{Beispiele:}\\
            der in Wien geborene Musiker : the musician born in Vienna\\
            die gestern abgeschickte E-Mail : the email sent yesterday\\
            die seit Jahren steigenden Kosten : the costs rising for years\\
            die von vielen Experten kritisierte Entscheidung : the decision criticized by many experts
    
    
    \section{Kasus {\small (Cases)}}
    
    \subsection{Nominativ, Akkusativ, Dativ und Genitiv {\small (Nominative, Accusative, Dative, and Genitive)}}
        
        \paragraph{Nominativ {\small (Nominative)}}
            Der Nominativ steht meist beim Subjekt oder nach bestimmten Verben wie \textit{sein}.\\
            The nominative usually appears with the subject or after certain verbs such as \textit{sein}.
            
            \textbf{Beispiele:}\\
            Der Student liest. : The student is reading.\\
            Die Frau ist Ärztin. : The woman is a doctor.\\
            Das Wetter ist schön. : The weather is nice.\\
            Mein Bruder kommt später. : My brother is coming later.
        
        \paragraph{Akkusativ {\small (Accusative)}}
            Der Akkusativ steht oft beim direkten Objekt und nach bestimmten Präpositionen.\\
            The accusative often appears with the direct object and after certain prepositions.
            
            \textbf{Beispiele:}\\
            Ich sehe den Mann. : I see the man.\\
            Sie kauft eine Tasche. : She buys a bag.\\
            Wir gehen durch den Park. : We walk through the park.\\
            Er besucht seinen Freund. : He visits his friend.
        
        \paragraph{Temporaler Akkusativ {\small (Temporal Accusative)}}
            Der temporale Akkusativ drückt Zeitdauer oder Zeitpunkt ohne Präposition aus.\\
            The temporal accusative expresses duration or time without a preposition.
            
            \textbf{Beispiele:}\\
            Ich warte den ganzen Tag. : I am waiting the whole day.\\
            Er bleibt eine Woche in Wien. : He stays in Vienna for a week.\\
            Wir arbeiten jeden Abend. : We work every evening.\\
            Sie schlief die ganze Nacht nicht. : She did not sleep all night.
        
        \paragraph{Dativ {\small (Dative)}}
            Der Dativ steht oft beim indirekten Objekt oder nach bestimmten Präpositionen und Verben.\\
            The dative often appears with the indirect object or after certain prepositions and verbs.
            
            \textbf{Beispiele:}\\
            Ich gebe dem Kind einen Apfel. : I give the child an apple.\\
            Sie hilft ihrem Vater. : She helps her father.\\
            Er wohnt bei seinen Eltern. : He lives with his parents.\\
            Wir danken Ihnen für Ihre Hilfe. : We thank you for your help.
        
        \paragraph{Freier Dativ {\small (Free Dative)}}
            Der freie Dativ drückt Beteiligung, Interesse oder Betroffenheit aus.\\
            The free dative expresses involvement, interest, or affectedness.
            
            \textbf{Beispiele:}\\
            Mach mir bitte die Tür zu. : Please close the door for me.\\
            Er ist mir zu laut. : He is too loud for me.\\
            Der Hund ist mir weggelaufen. : My dog ran away on me.\\
            Schreib mir bald. : Write to me soon.
        
        \paragraph{Genitiv {\small (Genitive)}}
            Der Genitiv drückt oft Zugehörigkeit oder nähere Bestimmung aus und erscheint auch nach bestimmten Präpositionen.\\
            The genitive often expresses possession or specification and also appears after certain prepositions.
            
            \textbf{Beispiele:}\\
            Das Auto des Arztes ist neu. : The doctor's car is new.\\
            Wegen des Regens bleiben wir zu Hause. : Because of the rain we stay at home.\\
            Trotz des Problems machten wir weiter. : Despite the problem we continued.\\
            Die Bedeutung dieser Frage ist groß. : The importance of this question is great.
    
    
    \section{Präpositionen {\small (Prepositions)}}
    
    \subsection{Präpositionen nach Kasus und Bedeutung {\small (Prepositions by Case and Meaning)}}
        
        \paragraph{Präpositionen mit Dativ {\small (Prepositions with Dative)}}
            Einige Präpositionen verlangen immer den Dativ.\\
            Some prepositions always require the dative.
            
            \textbf{Beispiele:}\\
            Ich wohne bei meinen Eltern. : I live with my parents.\\
            Sie fährt mit dem Bus. : She goes by bus.\\
            Er kommt aus der Schweiz. : He comes from Switzerland.\\
            Wir sprechen seit dem Morgen darüber. : We have been talking about it since the morning.
        
        \paragraph{Präpositionen mit Akkusativ {\small (Prepositions with Accusative)}}
            Andere Präpositionen stehen immer mit dem Akkusativ.\\
            Other prepositions always take the accusative.
            
            \textbf{Beispiele:}\\
            Wir gehen durch den Wald. : We are walking through the forest.\\
            Er arbeitet für einen Arzt. : He works for a doctor.\\
            Sie geht ohne ihren Bruder. : She goes without her brother.\\
            Ich komme gegen acht Uhr. : I am coming around eight o'clock.
        
        \paragraph{Präpositionen mit Dativ und Akkusativ {\small (Two-Way Prepositions)}}
            Wechselpräpositionen stehen mit Dativ bei Ort und mit Akkusativ bei Richtung.\\
            Two-way prepositions take the dative for location and the accusative for direction.
            
            \textbf{Beispiele:}\\
            Das Buch liegt auf dem Tisch. : The book is lying on the table.\\
            Ich lege das Buch auf den Tisch. : I put the book on the table.\\
            Sie sitzt in der Küche. : She is sitting in the kitchen.\\
            Er geht in die Küche. : He goes into the kitchen.
        
        \paragraph{Präpositionen mit Genitiv {\small (Prepositions with Genitive)}}
            Einige Präpositionen verlangen den Genitiv, besonders in formeller Sprache.\\
            Some prepositions require the genitive, especially in formal language.
            
            \textbf{Beispiele:}\\
            Während des Gesprächs blieb er ruhig. : During the conversation he remained calm.\\
            Trotz des Wetters gingen wir spazieren. : Despite the weather we went for a walk.\\
            Wegen eines Fehlers wurde alles verschoben. : Because of an error everything was postponed.\\
            Innerhalb eines Jahres hat sich viel verändert. : Within one year a lot changed.
        
        \paragraph{Lokale Präpositionen {\small (Local Prepositions)}}
            Lokale Präpositionen beschreiben Ort oder Richtung.\\
            Local prepositions describe place or direction.
            
            \textbf{Beispiele:}\\
            Das Bild hängt an der Wand. : The picture is hanging on the wall.\\
            Er geht an die Wand. : He goes to the wall.\\
            Sie steht vor dem Haus. : She is standing in front of the house.\\
            Wir laufen hinter das Gebäude. : We run behind the building.
        
        \paragraph{Temporale Präpositionen {\small (Temporal Prepositions)}}
            Temporale Präpositionen ordnen Ereignisse zeitlich ein.\\
            Temporal prepositions place events in time.
            
            \textbf{Beispiele:}\\
            Am Montag habe ich keine Zeit. : On Monday I have no time.\\
            Im Winter bleibe ich lieber zu Hause. : In winter I prefer to stay at home.\\
            Seit einem Jahr lernt er Deutsch. : He has been learning German for a year.\\
            Nach dem Essen trinken wir Tee. : After the meal we drink tea.
        
        \paragraph{Kausale und modale Präpositionen {\small (Causal and Modal Prepositions)}}
            Diese Präpositionen drücken Grund oder Art und Weise aus.\\
            These prepositions express reason or manner.
            
            \textbf{Beispiele:}\\
            Wegen der Krankheit blieb sie zu Hause. : Because of the illness she stayed at home.\\
            Dank deiner Hilfe war alles einfacher. : Thanks to your help everything was easier.\\
            Mit großer Freude nahm er das Angebot an. : He accepted the offer with great joy.\\
            Ohne viel zu sagen, ging er weg. : Without saying much, he left.
    
    
    \section{Satzbau {\small (Sentence Structure)}}
    
    \subsection{Wortstellung und Satzarten {\small (Word Order and Sentence Types)}}
        
        \paragraph{Aussagesatz und Verbzweitstellung {\small (Declarative Sentences and Verb-Second Position)}}
            Im Hauptsatz steht das konjugierte Verb normalerweise an zweiter Stelle.\\
            In the main clause, the conjugated verb normally stands in second position.
            
            \textbf{Beispiele:}\\
            Ich gehe heute ins Kino. : I am going to the cinema today.\\
            Heute gehe ich ins Kino. : Today I am going to the cinema.\\
            Morgen arbeitet er zu Hause. : Tomorrow he is working from home.\\
            Seit gestern wohnt sie in Wien. : Since yesterday she has been living in Vienna.
        
        \paragraph{Fragesatz {\small (Interrogative Sentence)}}
            Im Entscheidungsfragesatz steht das Verb an erster Stelle, bei W-Fragen nach dem Fragewort.\\
            In yes-no questions, the verb stands first; in wh-questions, it follows the question word.
            
            \textbf{Beispiele:}\\
            Kommst du morgen? : Are you coming tomorrow?\\
            Hast du das gesehen? : Did you see that?\\
            Warum bist du so müde? : Why are you so tired?\\
            Wann beginnt der Kurs? : When does the course begin?
        
        \paragraph{Satzklammer {\small (Sentence Bracket)}}
            Viele deutsche Sätze haben eine Satzklammer: ein Verbteil steht früh, der andere am Ende.\\
            Many German sentences have a sentence bracket: one verb part appears early, the other at the end.
            
            \textbf{Beispiele:}\\
            Ich habe das Buch gestern gekauft. : I bought the book yesterday.\\
            Er wird morgen ankommen. : He will arrive tomorrow.\\
            Sie möchte heute nicht ausgehen. : She does not want to go out today.\\
            Wir haben lange darüber nachgedacht. : We thought about it for a long time.
        
        \paragraph{Stellung von Satzgliedern {\small (Position of Sentence Elements)}}
            Die Reihenfolge im Mittelfeld ist flexibel, aber Informationsstruktur und Betonung spielen eine große Rolle.\\
            The order in the middle field is flexible, but information structure and emphasis play a major role.
            
            \textbf{Beispiele:}\\
            Ich habe ihm das Buch gegeben. : I gave him the book.\\
            Ich habe das Buch ihm gegeben. : I gave the book to him.\\
            Heute habe ich ihm das Buch gegeben. : Today I gave him the book.\\
            Wahrscheinlich hat er es schon gelesen. : He has probably already read it.
        
        \paragraph{Negation mit \textit{nicht} und \textit{kein} {\small (Negation with \textit{nicht} and \textit{kein})}}
            \textit{Nicht} negiert oft Verben, Adjektive oder Satzteile; \textit{kein} negiert Nomen.\\
            \textit{Nicht} often negates verbs, adjectives, or parts of the sentence; \textit{kein} negates nouns.
            
            \textbf{Beispiele:}\\
            Ich komme heute nicht. : I am not coming today.\\
            Das ist nicht einfach. : That is not easy.\\
            Ich habe kein Geld. : I have no money.\\
            Sie trinkt keinen Kaffee. : She drinks no coffee.
        
        \paragraph{Aufforderungssatz {\small (Imperative Sentence)}}
            Der Imperativ drückt Bitte, Aufforderung, Rat oder Befehl aus.\\
            The imperative expresses a request, instruction, advice, or command.
            
            \textbf{Beispiele:}\\
            Komm bitte herein. : Please come in.\\
            Seien Sie vorsichtig. : Be careful.\\
            Nehmt eure Bücher heraus. : Take out your books.\\
            Lass mich in Ruhe. : Leave me alone.
    
    
    \section{Nebensätze und Satzverbindungen {\small (Subordinate Clauses and Sentence Connections)}}
    
    \subsection{Grundtypen von Nebensätzen {\small (Basic Types of Subordinate Clauses)}}
        
        \paragraph{\textit{dass}-Sätze {\small (\textit{that}-clauses)}}
            \textit{Dass}-Sätze stehen oft nach Verben des Sagens, Denkens, Wissens oder Fühlens.\\
            \textit{Dass}-clauses often follow verbs of saying, thinking, knowing, or feeling.
            
            \textbf{Beispiele:}\\
            Ich glaube, dass er recht hat. : I think that he is right.\\
            Sie weiß, dass wir kommen. : She knows that we are coming.\\
            Er sagt, dass er keine Zeit hat. : He says that he has no time.\\
            Wir hoffen, dass alles gut wird. : We hope that everything will turn out well.
        
        \paragraph{\textit{ob}-Sätze {\small (\textit{whether}-clauses)}}
            \textit{Ob}-Sätze drücken eine indirekte Ja-Nein-Frage aus.\\
            \textit{Ob}-clauses express an indirect yes-no question.
            
            \textbf{Beispiele:}\\
            Ich weiß nicht, ob er kommt. : I do not know whether he is coming.\\
            Sie fragt, ob ich Zeit habe. : She asks whether I have time.\\
            Wir müssen klären, ob das möglich ist. : We need to clarify whether that is possible.\\
            Er überlegt, ob er umziehen soll. : He is considering whether he should move.
        
        \paragraph{Indirekte Fragesätze {\small (Indirect Questions)}}
            Indirekte Fragesätze beginnen mit einem Fragewort und haben Verbendstellung.\\
            Indirect questions begin with a question word and have verb-final order.
            
            \textbf{Beispiele:}\\
            Ich weiß nicht, warum er so spät kommt. : I do not know why he is coming so late.\\
            Sie fragt, wann der Kurs beginnt. : She asks when the course starts.\\
            Er erklärt, wie das Gerät funktioniert. : He explains how the device works.\\
            Wir wissen nicht, wo sie wohnt. : We do not know where she lives.
        
        \paragraph{Relativsätze {\small (Relative Clauses)}}
            Relativsätze geben zusätzliche Informationen zu einem Nomen.\\
            Relative clauses provide additional information about a noun.
            
            \textbf{Beispiele:}\\
            Das ist der Mann, der hier arbeitet. : That is the man who works here.\\
            Ich kenne die Frau, die dort steht. : I know the woman who is standing there.\\
            Das ist das Buch, das du suchst. : That is the book you are looking for.\\
            Wir besuchen Freunde, die in Graz wohnen. : We are visiting friends who live in Graz.
        
        \paragraph{Relativsätze mit Präposition {\small (Relative Clauses with Prepositions)}}
            Wenn das Bezugswort eine Präposition verlangt, erscheint sie im Relativsatz vor dem Relativpronomen.\\
            If the reference word requires a preposition, it appears before the relative pronoun in the relative clause.
            
            \textbf{Beispiele:}\\
            Das ist die Frau, mit der ich gesprochen habe. : That is the woman I spoke with.\\
            Das ist das Thema, über das wir diskutieren. : That is the topic we are discussing.\\
            Hier ist der Stuhl, auf dem ich saß. : Here is the chair on which I sat.\\
            Das ist der Grund, aus dem er gegangen ist. : That is the reason why he left.
        
        \paragraph{Relativsätze im Genitiv {\small (Genitive Relative Clauses)}}
            Mit \textit{dessen} und \textit{deren} drückt man Besitz oder Zugehörigkeit im Relativsatz aus.\\
            With \textit{dessen} and \textit{deren}, one expresses possession or belonging in a relative clause.
            
            \textbf{Beispiele:}\\
            Das ist der Mann, dessen Auto gestohlen wurde. : That is the man whose car was stolen.\\
            Sie ist die Frau, deren Bruder Arzt ist. : She is the woman whose brother is a doctor.\\
            Das ist das Kind, dessen Mutter Lehrerin ist. : That is the child whose mother is a teacher.\\
            Wir kennen den Autor, dessen Buch berühmt wurde. : We know the author whose book became famous.
        
        \paragraph{Infinitivsätze mit \textit{zu} {\small (Infinitive Clauses with \textit{zu})}}
            Infinitivsätze mit \textit{zu} verdichten Informationen und treten oft nach bestimmten Verben, Adjektiven und Nomen auf.\\
            Infinitive clauses with \textit{zu} condense information and often occur after certain verbs, adjectives, and nouns.
            
            \textbf{Beispiele:}\\
            Ich versuche, pünktlich zu kommen. : I try to come on time.\\
            Er hofft, die Prüfung zu bestehen. : He hopes to pass the exam.\\
            Es ist wichtig, genug zu schlafen. : It is important to sleep enough.\\
            Sie hat keine Lust, heute zu kochen. : She does not feel like cooking today.
        
        \paragraph{Verkürzte Relativsätze und Partizipialkonstruktionen {\small (Reduced Relative Clauses and Participial Constructions)}}
            In schriftlichen Texten werden Relativsätze oft durch Partizipialkonstruktionen verkürzt.\\
            In written texts, relative clauses are often reduced by participial constructions.
            
            \textbf{Beispiele:}\\
            die gestern angekommenen Gäste : the guests who arrived yesterday\\
            die im Zentrum liegende Wohnung : the apartment located in the center\\
            die von vielen kritisierte Entscheidung : the decision criticized by many\\
            die ständig steigenden Preise : the constantly rising prices
    
    \subsection{Adverbiale Nebensätze {\small (Adverbial Clauses)}}
    
    \paragraph{Kausale Nebensätze {\small (Causal Clauses)}}
        Kausale Nebensätze geben einen Grund an.\\
        Causal clauses express a reason.
        
        \textbf{Beispiele:}\\
        Ich bleibe zu Hause, weil ich krank bin. : I stay at home because I am ill.\\
        Da es regnet, nehmen wir ein Taxi. : Since it is raining, we take a taxi.\\
        Weil er müde war, ging er früh ins Bett. : Because he was tired, he went to bed early.\\
        Da ich keine Zeit hatte, sagte ich ab. : Since I had no time, I canceled.
    
    \paragraph{Konzessive Nebensätze {\small (Concessive Clauses)}}
        Konzessive Nebensätze drücken einen Gegensatz oder eine Einschränkung aus.\\
        Concessive clauses express contrast or concession.
        
        \textbf{Beispiele:}\\
        Obwohl es spät war, arbeitete sie weiter. : Although it was late, she kept working.\\
        Auch wenn du recht hast, musst du ruhig bleiben. : Even if you are right, you must stay calm.\\
        Obgleich er krank war, kam er zur Arbeit. : Although he was ill, he came to work.\\
        Selbst wenn es regnet, gehen wir spazieren. : Even if it rains, we will go for a walk.
    
    \paragraph{Konditionale Nebensätze {\small (Conditional Clauses)}}
        Konditionale Nebensätze nennen eine Bedingung.\\
        Conditional clauses state a condition.
        
        \textbf{Beispiele:}\\
        Wenn ich Zeit habe, komme ich vorbei. : If I have time, I will come by.\\
        Falls du Hilfe brauchst, ruf mich an. : If you need help, call me.\\
        Sofern alles klappt, reisen wir morgen ab. : Provided everything works out, we leave tomorrow.\\
        Wenn er fragt, sage ich die Wahrheit. : If he asks, I will tell the truth.
    
    \paragraph{Modale Nebensätze {\small (Modal Clauses)}}
        Modale Nebensätze beschreiben die Art und Weise eines Vorgangs.\\
        Modal clauses describe the manner in which something happens.
        
        \textbf{Beispiele:}\\
        Er sah mich an, als ob er mich kennen würde. : He looked at me as if he knew me.\\
        Sie spricht, indem sie langsam und deutlich formuliert. : She speaks by formulating slowly and clearly.\\
        Ohne dass er etwas sagte, ging er weg. : Without saying anything, he left.\\
        Anstatt dass sie wartet, beginnt sie sofort. : Instead of waiting, she starts immediately.
    
    \paragraph{Temporale Nebensätze {\small (Temporal Clauses)}}
        Temporale Nebensätze ordnen Handlungen zeitlich.\\
        Temporal clauses arrange actions in time.
        
        \textbf{Beispiele:}\\
        Als ich klein war, wohnte ich auf dem Land. : When I was little, I lived in the countryside.\\
        Wenn ich nach Hause komme, koche ich. : When I come home, I cook.\\
        Nachdem wir gegessen hatten, gingen wir spazieren. : After we had eaten, we went for a walk.\\
        Bevor du gehst, ruf mich bitte an. : Before you go, please call me.
    
    \paragraph{\textit{seitdem}, \textit{bis}, \textit{sobald}, \textit{solange}, \textit{während} {\small (Specific Temporal Connectors)}}
        Diese Konjunktionen geben genauere zeitliche Beziehungen an.\\
        These conjunctions express more specific temporal relations.
        
        \textbf{Beispiele:}\\
        Seitdem ich hier wohne, bin ich ruhiger. : Since I have lived here, I have become calmer.\\
        Warte hier, bis ich zurückkomme. : Wait here until I come back.\\
        Sobald er ankommt, beginnen wir. : As soon as he arrives, we begin.\\
        Solange du lernst, wirst du Fortschritte machen. : As long as you study, you will make progress.\\
        Während ich arbeite, höre ich Musik. : While I work, I listen to music.
    
    \paragraph{Konsekutive Nebensätze {\small (Consecutive Clauses)}}
        Konsekutive Nebensätze zeigen eine Folge oder ein Ergebnis.\\
        Consecutive clauses show a consequence or result.
        
        \textbf{Beispiele:}\\
        Er war so müde, dass er sofort einschlief. : He was so tired that he fell asleep immediately.\\
        Sie sprach so leise, dass niemand sie verstand. : She spoke so quietly that nobody understood her.\\
        Es war zu spät, als dass wir noch hätten beginnen können. : It was too late for us to still begin.\\
        Der Film war so spannend, dass wir die Zeit vergaßen. : The film was so exciting that we forgot the time.
    
    \paragraph{Adversative Nebensätze und Verbindungen {\small (Adversative Clauses and Connections)}}
        Adversative Verbindungen stellen Gegensätze oder Unterschiede dar.\\
        Adversative structures present contrast or difference.
        
        \textbf{Beispiele:}\\
        Während er gern reist, bleibt sie lieber zu Hause. : While he likes to travel, she prefers to stay at home.\\
        Wohingegen sie Ruhe sucht, liebt er das Stadtleben. : Whereas she seeks peace, he loves city life.\\
        Er ist freundlich, aber sehr direkt. : He is friendly, but very direct.\\
        Sie wollte kommen, jedoch war sie krank. : She wanted to come, however she was ill.
    
    \paragraph{Finalsätze {\small (Purpose Clauses)}}
        Finalsätze drücken Zweck oder Ziel aus.\\
        Purpose clauses express aim or purpose.
        
        \textbf{Beispiele:}\\
        Ich lerne Deutsch, damit ich in Österreich arbeiten kann. : I am learning German so that I can work in Austria.\\
        Er spart Geld, um ein Auto zu kaufen. : He is saving money in order to buy a car.\\
        Sie wiederholt alles, damit sie nichts vergisst. : She repeats everything so that she forgets nothing.\\
        Wir fahren früh los, um pünktlich anzukommen. : We leave early in order to arrive on time.
    
    \subsection{Konjunktionen und Konnektoren {\small (Conjunctions and Connectors)}}
    
    \paragraph{Nebenordnende und unterordnende Konjunktionen {\small (Coordinating and Subordinating Conjunctions)}}
        Nebenordnende Konjunktionen verbinden gleichrangige Satzteile, unterordnende leiten Nebensätze ein.\\
        Coordinating conjunctions connect elements of equal rank, subordinating ones introduce subordinate clauses.
        
        \textbf{Beispiele:}\\
        Ich komme, aber er bleibt zu Hause. : I am coming, but he is staying at home.\\
        Sie arbeitet viel und lernt zusätzlich Deutsch. : She works a lot and also learns German.\\
        Er geht nicht, weil er krank ist. : He is not going because he is ill.\\
        Obwohl es regnet, fahren wir los. : Although it is raining, we set off.
    
    \paragraph{Satzverbindende Konnektoren {\small (Sentence-Linking Connectors)}}
        Konnektoren strukturieren Texte und machen Beziehungen zwischen Aussagen klar.\\
        Connectors structure texts and clarify relationships between statements.
        
        \textbf{Beispiele:}\\
        Es regnete. Deshalb blieben wir zu Hause. : It rained. Therefore we stayed at home.\\
        Er war krank. Trotzdem kam er zur Arbeit. : He was ill. Nevertheless he came to work.\\
        Die Idee ist gut. Allerdings ist sie teuer. : The idea is good. However, it is expensive.\\
        Einerseits ist die Lage schwierig, andererseits gibt es Hoffnung. : On the one hand the situation is difficult, on the other hand there is hope.
    
    \paragraph{Mehrteilige Konnektoren {\small (Multi-Part Connectors)}}
        Mehrteilige Konnektoren sind besonders wichtig für klare Argumentation und gehobenen Stil.\\
        Multi-part connectors are especially important for clear argumentation and a more advanced style.
        
        \textbf{Beispiele:}\\
        Sowohl der Lehrer als auch die Studenten waren zufrieden. : Both the teacher and the students were satisfied.\\
        Er ist zwar jung, aber sehr erfahren. : He is indeed young, but very experienced.\\
        Weder der Arzt noch die Apotheke war erreichbar. : Neither the doctor nor the pharmacy was reachable.\\
        Nicht nur die Grammatik, sondern auch der Wortschatz ist wichtig. : Not only grammar but also vocabulary is important.
    
    
    \section{Passiv {\small (Passive Voice)}}
    
    \subsection{Vorgangspassiv und Zustandspassiv {\small (Process Passive and State Passive)}}
        
        \paragraph{Vorgangspassiv {\small (Process Passive)}}
            Das Vorgangspassiv beschreibt eine Handlung oder einen Vorgang. Es wird mit \textit{werden} und Partizip II gebildet.\\
            The process passive describes an action or process. It is formed with \textit{werden} and the past participle.
            
            \textbf{Beispiele:}\\
            Das Haus wird gebaut. : The house is being built.\\
            Die E-Mail wurde gestern geschickt. : The email was sent yesterday.\\
            Die Straße wird repariert. : The street is being repaired.\\
            Die Regeln werden erklärt. : The rules are being explained.
        
        \paragraph{Vorgangspassiv mit Modalverben {\small (Process Passive with Modal Verbs)}}
            Auch im Passiv können Modalverben stehen.\\
            Modal verbs can also be used in the passive.
            
            \textbf{Beispiele:}\\
            Die Aufgabe muss heute erledigt werden. : The task must be completed today.\\
            Das Problem kann schnell gelöst werden. : The problem can be solved quickly.\\
            Der Antrag soll morgen bearbeitet werden. : The application is supposed to be processed tomorrow.\\
            Die Daten dürfen nicht weitergegeben werden. : The data may not be passed on.
        
        \paragraph{Zustandspassiv {\small (State Passive)}}
            Das Zustandspassiv beschreibt einen Zustand als Ergebnis einer Handlung. Es wird mit \textit{sein} und Partizip II gebildet.\\
            The state passive describes a state as the result of an action. It is formed with \textit{sein} and the past participle.
            
            \textbf{Beispiele:}\\
            Die Tür ist geschlossen. : The door is closed.\\
            Das Fenster ist geöffnet. : The window is open.\\
            Der Brief ist unterschrieben. : The letter is signed.\\
            Die Aufgabe ist erledigt. : The task is completed.
        
        \paragraph{Passiversatzformen {\small (Passive Alternatives)}}
            Im Deutschen gibt es mehrere aktive oder halbpassive Alternativen zum Passiv.\\
            In German there are several active or semi-passive alternatives to the passive.
            
            \textbf{Beispiele:}\\
            Man baut das Haus neu. : They are rebuilding the house.\\
            Das Haus lässt sich leicht renovieren. : The house can be renovated easily.\\
            Der Bericht ist noch zu schreiben. : The report still has to be written.\\
            Ich bekam die Unterlagen zugeschickt. : I was sent the documents.
    
    
    \section{Modalität {\small (Modality)}}
    
    \subsection{Objektive und subjektive Modalität {\small (Objective and Subjective Modality)}}
        
        \paragraph{Objektive Bedeutung der Modalverben {\small (Objective Meaning of Modal Verbs)}}
            Hier geht es um reale Fähigkeit, Erlaubnis, Notwendigkeit oder Verpflichtung.\\
            Here, the focus is on real ability, permission, necessity, or obligation.
            
            \textbf{Beispiele:}\\
            Du darfst jetzt gehen. : You may go now.\\
            Ich muss morgen arbeiten. : I have to work tomorrow.\\
            Sie kann sehr schnell lesen. : She can read very quickly.\\
            Wir sollen pünktlich sein. : We are supposed to be on time.
        
        \paragraph{Subjektive Bedeutung der Modalverben {\small (Subjective Meaning of Modal Verbs)}}
            Subjektive Modalität drückt Vermutung, Unsicherheit oder Schlussfolgerung aus.\\
            Subjective modality expresses assumption, uncertainty, or inference.
            
            \textbf{Beispiele:}\\
            Er muss schon zu Hause sein. : He must already be at home.\\
            Sie dürfte etwa dreißig sein. : She is probably about thirty.\\
            Das kann nicht stimmen. : That cannot be right.\\
            Er will krank sein. : He claims to be ill.
        
        \paragraph{Satzadverbien der Einschätzung {\small (Sentence Adverbs of Assessment)}}
            Mit Satzadverbien kann man Sicherheit, Unsicherheit oder Bewertung ausdrücken.\\
            With sentence adverbs one can express certainty, uncertainty, or evaluation.
            
            \textbf{Beispiele:}\\
            Wahrscheinlich kommt er später. : He will probably come later.\\
            Offenbar wusste niemand Bescheid. : Apparently nobody knew about it.\\
            Möglicherweise hat sie den Zug verpasst. : She may have missed the train.\\
            Jedenfalls müssen wir handeln. : In any case, we must act.
    
    
    \section{Umformungen und Stil {\small (Transformations and Style)}}
    
    \subsection{Stilistische und syntaktische Umformungen {\small (Stylistic and Syntactic Transformations)}}
        
        \paragraph{Nominalisierung {\small (Nominalization)}}
            Nominalisierungen verdichten Informationen und sind typisch für wissenschaftliche oder formelle Sprache.\\
            Nominalizations condense information and are typical of academic or formal language.
            
            \textbf{Beispiele:}\\
            die Entscheidung des Gerichts : the decision of the court\\
            nach der Ankunft : after the arrival\\
            beim Lesen des Textes : while reading the text\\
            die Verbesserung der Situation : the improvement of the situation
        
        \paragraph{Verbalisierung {\small (Verbalization)}}
            Bei der Verbalisierung werden nominale Ausdrücke in verbalere Strukturen umgewandelt.\\
            In verbalization, nominal expressions are transformed into more verbal structures.
            
            \textbf{Beispiele:}\\
            Nach der Ankunft begann die Sitzung. : After the arrival, the meeting began.\\
            Nachdem wir angekommen waren, begann die Sitzung. : After we had arrived, the meeting began.\\
            Die Entscheidung des Teams überraschte alle. : The team's decision surprised everyone.\\
            Dass das Team so entschied, überraschte alle. : That the team decided so surprised everyone.
        
        \paragraph{Umwandlung von Präpositionalgruppen in Nebensätze {\small (Transformation of Prepositional Phrases into Subordinate Clauses)}}
            Präpositionalgruppen können oft durch Nebensätze ersetzt werden.\\
            Prepositional phrases can often be replaced by subordinate clauses.
            
            \textbf{Beispiele:}\\
            Wegen des Regens blieben wir zu Hause. : Because of the rain we stayed at home.\\
            Weil es regnete, blieben wir zu Hause. : Because it was raining, we stayed at home.\\
            Trotz seiner Müdigkeit arbeitete er weiter. : Despite his tiredness he kept working.\\
            Obwohl er müde war, arbeitete er weiter. : Although he was tired, he kept working.
        
        \paragraph{Umformung von Nebensätzen in Infinitivkonstruktionen {\small (Transforming Subordinate Clauses into Infinitive Constructions)}}
            Wenn das Subjekt gleich bleibt, kann ein Nebensatz oft in einen Infinitivsatz umgeformt werden.\\
            If the subject remains the same, a subordinate clause can often be transformed into an infinitive clause.
            
            \textbf{Beispiele:}\\
            Er ging früh, um pünktlich zu sein. : He left early in order to be on time.\\
            Er ging früh, damit er pünktlich war. : He left early so that he would be on time.\\
            Sie öffnete das Fenster, ohne etwas zu sagen. : She opened the window without saying anything.\\
            Anstatt zu warten, begann sie sofort. : Instead of waiting, she started immediately.
        
        \paragraph{Umformung von Aktiv und Passiv {\small (Transformation between Active and Passive)}}
            Aktiv und Passiv setzen unterschiedliche Schwerpunkte im Satz.\\
            Active and passive voice place emphasis differently in a sentence.
            
            \textbf{Beispiele:}\\
            Der Arzt untersucht den Patienten. : The doctor examines the patient.\\
            Der Patient wird vom Arzt untersucht. : The patient is examined by the doctor.\\
            Man diskutiert das Problem intensiv. : People are discussing the problem intensively.\\
            Das Problem wird intensiv diskutiert. : The problem is being discussed intensively.
        
        \paragraph{Indirekte Rede {\small (Indirect Speech)}}
            Indirekte Rede gibt Aussagen anderer Personen wieder, oft im Konjunktiv I.\\
            Indirect speech reports what other people have said, often using Subjunctive I.
            
            \textbf{Beispiele:}\\
            Er sagte, er habe keine Zeit. : He said that he had no time.\\
            Sie erklärte, sie werde später kommen. : She explained that she would come later.\\
            Der Zeuge meinte, er sei zu Hause gewesen. : The witness said that he had been at home.\\
            Die Presse berichtet, die Lage habe sich beruhigt. : The press reports that the situation has calmed down.
        
        \paragraph{Berichtsstil und Distanzierung {\small (Report Style and Distancing)}}
            Mit bestimmten Formen schafft man Distanz und vermeidet, eine Aussage als eigene Meinung darzustellen.\\
            With certain forms, one creates distance and avoids presenting a statement as one's own opinion.
            
            \textbf{Beispiele:}\\
            Es wird berichtet, dass die Zahlen steigen. : It is reported that the numbers are increasing.\\
            Angeblich war niemand informiert. : Allegedly nobody was informed.\\
            Dem Bericht zufolge hat sich die Lage verbessert. : According to the report, the situation has improved.\\
            Wie es heißt, soll das Projekt gestoppt werden. : According to reports, the project is to be stopped.
    
    
    \section{Weitere zentrale Bereiche {\small (Further Central Areas)}}
    
    \subsection{Wortbildung {\small (Word Formation)}}
        
        \paragraph{Präfixe, Suffixe und Wortfamilien {\small (Prefixes, Suffixes, and Word Families)}}
            Wortbildung hilft beim Verstehen neuer Wörter und beim Aufbau eines präzisen Wortschatzes.\\
            Word formation helps in understanding new words and building precise vocabulary.
            
            \textbf{Beispiele:}\\
            möglich -- unmöglich : possible -- impossible\\
            Freund -- freundlich -- Freundlichkeit : friend -- friendly -- friendliness\\
            fahren -- Abfahrt -- Fahrer -- Fahrzeug : drive -- departure -- driver -- vehicle\\
            sicher -- unsicher -- Sicherheit : secure -- insecure -- security
    
    \subsection{Wortstellung im erweiterten Satz {\small (Word Order in Extended Sentences)}}
    
    \paragraph{Informationsstruktur und Schwerpunktsetzung {\small (Information Structure and Focus)}}
        Die Stellung einzelner Satzteile kann den Schwerpunkt der Aussage verändern.\\
        The position of individual sentence parts can change the focus of the statement.
        
        \textbf{Beispiele:}\\
        Heute habe ich den Bericht geschrieben. : Today I wrote the report.\\
        Den Bericht habe ich heute geschrieben. : It was the report that I wrote today.\\
        Ich habe heute den Bericht geschrieben. : I wrote the report today.\\
        Geschrieben habe ich den Bericht heute noch nicht. : Written the report I have not yet today.
    
    \subsection{Ellipsen und verkürzte Strukturen {\small (Ellipses and Reduced Structures)}}
    
    \paragraph{Auslassungen in gesprochener und schriftlicher Sprache {\small (Omissions in Spoken and Written Language)}}
        Nicht immer wird ein vollständiger Satz verwendet. In Gesprächen und in bestimmten Textsorten sind Ellipsen häufig.\\
        A full sentence is not always used. In conversations and certain text types, ellipses are frequent.
        
        \textbf{Beispiele:}\\
        Noch Fragen? : Any more questions?\\
        Alles klar. : Everything clear.\\
        Je früher, desto besser. : The sooner, the better.\\
        Wenn möglich, morgen. : If possible, tomorrow.
    
    \subsection{Stilregister und Ausdruck {\small (Registers and Expression)}}
    
    \paragraph{Alltagssprache, Standardsprache und formeller Stil {\small (Everyday Language, Standard Language, and Formal Style)}}
        Deutsch kann je nach Situation direkter, neutraler oder formeller klingen. Ein bewusster Wechsel des Registers ist wichtig.\\
        German can sound more direct, neutral, or formal depending on the situation. A conscious change of register is important.
        
        \textbf{Beispiele:}\\
        Kannst du mir helfen? : Can you help me?\\
        Könnten Sie mir bitte helfen? : Could you please help me?\\
        Ich möchte Sie höflich darum bitten, mir behilflich zu sein. : I would politely like to ask you to assist me.\\
        Das ist nicht gut. / Das ist problematisch. : That is not good. / That is problematic.

\end{document}